Выбирая стек технологий, мы прежде всего опирались на возможность создать уникальный проект, который позволит нам изучить разработку компьютерных игр. Мы выбрали писать игру на языке программирования С++, используя  SFML, так как в дальнейшем мы планируем развивать навыки рабоы с игровыми движками, а SFML "--- лучший выбор для начинающих программистов, которые выбирают разработку 2Dигр на С++.

\begin{table}[H]
\centering
\caption{Преимущества C++ при разработке игр}
\label{tab:cppDev}

\begin{tabular}{|p{0.3\textwidth}|p{0.3\textwidth}|p{0.35\textwidth}|}
\hline
\textbf{Характеристика} & \textbf{Преимущество C++} & \textbf{Практическое применение в играх} \\
\hline
Производительность & Близость к машинному коду, минимальные накладные расходы & Рендеринг сложных 3D-сцен, физические симуляции \\
\hline
Управление памятью & Полный контроль над выделением и освобождением ресурсов & Стабильный FPS, отсутствие фризов при сборке мусора \\
\hline
Кроссплатформенность & Компиляция для различных архитектур & Одновременный выпуск на ПК, консолях и мобильных устройствах \\
\hline
Экосистема & Множество библиотек и инструментов & Интеграция с физическими движками, системами рендеринга \\
\hline
Совместимость & Легкая интеграция с кодом на C и ассемблере & Оптимизация критических участков кода до уровня ассемблера \\
\hline
\end{tabular}
\end{table}

В нашем случае, именно C++ был тем языком программирования, который позволял создать качественный продукт. В таблице ~\ref{tab:cppDev} показано, какие преимущества есть в разработке компьютерных игр с нашим стеком технологий. На языке С++ было написано большое количество игр, которые известны сейчас почти каждому (Undertale, Terraria и другие). C++ занимает уникальное положение в мире разработки игр по целому ряду причин, которые делают его практически незаменимым для создания высокопроизводительных игр любого масштаба. 

Прежде всего, C++ предлагает беспрецедентную производительность. В отличие от языков с автоматическим управлением памятью, таких как Java или C\#, C++ даёт разработчикам прямой доступ к аппаратным ресурсам. Когда речь идёт о рендеринге сложных 3D"=сцен с миллионами полигонов или расчёте физики сотен объектов в реальном времени, каждая миллисекунда имеет значение.\cite{aleksC++}

Вторым ключевым преимуществом C++ является контроль над памятью. В играх, где требуется обрабатывать огромные объёмы данных в реальном времени, ручное управление памятью позволяет избежать непредсказуемых пауз, связанных с работой сборщика мусора. Разработчики могут точно определять, когда выделять и освобождать память, что критически важно для поддержания стабильного количества кадров в секунду.

Также еще одним преимуществом языка программирования С++ является наличие в нем объектно"=ориентированного подхода, который позволяет сделать код понятнее (например, классы для состояний объектов).\cite{OOPC++}

Разработка игр на C++ редко начинается с чистого листа. Современные разработчики опираются на богатую экосистему библиотек и фреймворков, которые значительно ускоряют процесс создания игр, предоставляя готовые решения для типичных задач геймдева. Умение выбрать и эффективно использовать эти инструменты часто определяет успех проекта не меньше, чем знание самого языка C++. Для разработки своей игры мы выбрали библиотеку SFML.

SFML (Simple and Fast Multimedia Library) "--- это свободная кроссплатформенная мультимедийная библиотека, написанная на C++ с объектно"=ориентированным подходом, служащая простым и быстрым интерфейсом для работы с графикой, звуком, вводом и сетью. Библиотека состоит из пяти модулей: System (время, потоки), Window (создание окон, обработка событий), Graphics (рисование спрайтов, текстур, шейдеров, поддержка OpenGL), Audio (звуки, музыка) и Network (TCP/UDP-соединения). Это позволяет создавать 2D"=игры, интерактивные приложения, визуализации данных или даже базовые 3D"=проекты с OpenGL без лишней сложности. 
